% !TeX program = lualatex
% !TeX encoding = utf8
% !TeX spellcheck = uk_UA
% !TeX root = ../main.tex

%=========================================================
\chapter{Плазма}\label{\currfilebase}\hypertarget{Plasma}{}
\graphicspath{{\currfilebase/Pictures}}
%=========================================================

%% --------------------------------------------------------
\section{Плазма}\label{sec:plasma}\hypertarget{plasma}{}
%% --------------------------------------------------------


	\emphx{Плазмою називається іонізований квазінейтральний газ, який займає настільки великий об'єм, що в ньому не відбувається скільки-небудь помітного порушення квазінейтральності через теплові флуктуації.} Квазінейтральність газу означає, що кількості позитивних і негативних зарядів у ньому \emphx{майже однакові}.

	      Оцінимо розміри області, в якій можуть відбуватися помітні порушення квазінейтральності, припускаючи для простоти, що заряди позитивних і негативних частинок однакові і рівні елементарному заряду $e$. Нехай така плазма заповнює простір між площинами $AB$ і $MN$ (рис.~\ref{pic:charge-separation}). Допустимо, що через теплові флуктуації негативні заряди змістилися вгору на відстань $l$. Тоді на границях плазми виникнуть макроскопічні заряди протилежних знаків з поверхневою густиною $\sigma = nle$, де $n$ --- концентрація частинок одного знака заряду. Напруженість електричного поля в плазмі буде $E = 4\pi\sigma = 4\pi nle$, а густина електричної енергії $E^2/8\pi = 2\pi(nle)^2$. Оскільки енергія електричного поля черпається з кінетичної енергії теплового руху частинок газу, величина $E^2/8\pi$ не може перевищувати $3nkT$. (На долю негативних частинок одиниці об'єму приходиться в середньому кінетична енергія $(3/2)kT$, і така ж енергія --- на долю позитивних.) Отже, якщо опустити числовий коефіцієнт 3, то повинно бути $2\pi(nle)^2 < nkT$, або $l < D$, де

\begin{equation}\label{eq:debye-length}
	D = \sqrt{\frac{kT}{2\pi ne^2}}.
\end{equation}

Величина $D$ називається \emphx{дебаївською довжиною} або \emphx{дебаївським радіусом}. (Більш точне визначення дається в задачі~1.)

Таким чином, \emphx{щоб плазма зберігала квазінейтральність, її лінійні розміри повинні набагато перевищувати дебаївський радіус $D$.} Тільки при дотриманні умови квазінейтральності плазма веде себе як \emphx{зв'язаний колектив заряджених частинок}, а не як проста сукупність невзаємодіючих частинок. Наприклад, якщо в плазмі є градієнт концентрації, то, незважаючи на різницю коефіцієнтів
дифузії, позитивні й негативні іони не дифундують з різними швидкостями, а завдяки квазінейтральності плазми переміщаються з однією й тією ж середньою швидкістю в одному й тому ж напрямку. Така дифузія називається \emphx{амбіполярною}. Фізична причина амбіполярного характеру дифузії плазми полягає в тому, що через різницю в коефіцієнтах дифузії та рухливості іонів у плазмі виникає деяке розділення електричних зарядів і пов'язане з ним електричне поле, яке сповільнює дифузію іонів одного знака й прискорює дифузію іонів протилежного знака. У результаті швидкості дифузії позитивних і негативних іонів вирівнюються. Відновлення порушеної квазінейтральності плазми аналогічне появі відновлювальних сил у пружних тілах при їхніх деформаціях. З цим пов'язана можливість різноманітних \emphx{колективних коливань} плазми, значно більш різноманітних, ніж у газах, що складаються з нейтральних частинок.

\begin{figure}[h!]\centering
	% \includegraphics[width=0.4\linewidth]{charge-separation}
	\caption{Розділення зарядів у плазмі}
	\label{pic:charge-separation}
\end{figure}


Залежно від ступеня іонізації $\alpha$ розрізняють \emphx{слабо іонізовану} (при $\alpha$ порядку долей процента), \emphx{помірно іонізовану} ($\alpha$ порядку декількох процентів) і \emphx{повністю іонізовану плазму}. У земних природних умовах плазма зустрічається досить рідко (наприклад, у каналі блискавки). У верхніх шарах атмосфери, які більшою мірою піддаються впливу іонізуючих факторів (ультрафіолетові та космічні промені), постійно присутня слабо іонізована плазма --- \emphx{іоносфера}, яка відбиває радіохвилі й робить можливим радіозв'язок на великих відстанях (порядку відстані між діаметрально протилежними точками земної кулі). У космічному просторі плазма являє собою найбільш поширений стан речовини. Сонце й гарячі зорі, які мають високі температури, складаються з повністю іонізованої плазми. Тому багато проблем астрофізики пов'язані з вивченням фізичних властивостей плазми. На ґрунті астрофізики виникла \emphx{магнітна гідродинаміка}, в якій плазма, що рухається в магнітних полях, розглядається як \emphx{суцільне рідке середовище}, що має \emphx{високу провідність}. Плазма утворюється в різних формах газового розряду, наприклад у позитивному стовпі тліючого розряду, а також у головному каналі іскрового розряду. Фізика плазми --- порівняно новий, швидко зростаючий розділ фізики, якому присвячені спеціальні курси. У курсі загальної фізики можна повідомити тільки деякі уривчасті відомості про плазму.

Оцінимо питому провідність $\lambda$ повністю іонізованої плазми, що складається з електронів і позитивно заряджених іонів із зарядом $Ze$ кожний. Рух іонів, зважаючи на їхні більші маси, можна не враховувати й вважати, що весь ток створюється рухом легких електронів. Величина $\lambda$ визначається зіткненнями електронів з іонами. Зіткнення електронів між собою на величину тока не впливають, оскільки при таких зіткненнях повний імпульс електронів не змінюється. Від цих зіткнень можна відволіктися. Між
іонами й електронами плазми діють \emphx{кулонівські сили притягання}. Це --- \emphx{далекодіючі сили}. Електрон порівняно рідко підходить до іона на такі малі відстані, щоб напрямок його руху змінився різко і мав характер стрибка. Набагато більше значення мають взаємодії електрона відразу з \emphx{дуже великою кількістю іонів}, при яких напрямок траєкторії електрона змінюється плавно й безперервно. Відхилення електрона на великі кути від першопочаткового напрямку руху відбуваються в результаті накопичення \emphx{малих відхилень} при взаємодії його з «далекими» іонами. Тому про зіткнення, довжину й час вільного пробігу можна говорити лише в умовному смислі. Умовимося називати часом вільного пробігу електрона проміжок часу $\tau$, протягом якого напрямок руху електрона змінюється на кут порядка $90^\circ$.

Для оцінки величини $\tau$ допустимо спочатку, що єдиний електрон рухається в полі позитивного іона із зарядом $Ze$. Якщо $v$ --- швидкість електрона на нескінченності, а $b$ --- прицільний параметр, то при проходженні мимо іона траєкторія електрона відхиляється на кут, який визначається формулою

\begin{equation*}
	\ctg\frac{\vartheta}{2} = \frac{mv^2b}{Ze^2},
\end{equation*}

де $m$ --- маса електрона (див.~т.~I, \S~58, де аналогічна формула отримана для рухів під дією гравітаційних сил). Прицільний параметр $b$, для якого $\vartheta = 90^\circ$, визначається виразом $b = Ze^2/(mv^2)$. Йому відповідає «ефективний поперечний переріз»

\begin{equation*}
	\sigma = \pi b^2 = \pi\left(\frac{Ze^2}{mv^2}\right)^2.
\end{equation*}

Врахування далеких взаємодій приводить до того ж результату, але збільшеному в $L$ разів:

\begin{equation*}
	\sigma = \pi L\left(\frac{Ze^2}{mv^2}\right)^2.
\end{equation*}

Коефіцієнт $L$ називається \emphx{кулонівським логарифмом}. Він майже не залежить від температури й густини плазми. Для плазми, що складається з повністю іонізованого дейтерію, при $kT \sim 10$~кеВ і концентрації електронів $n \sim 10^{12}$--$10^{15}$~см$^{-3}$ $L \approx 15$. Так як кожний позитивний іон містить $Z$ елементарних зарядів, то концентрація таких іонів буде $n/Z$, а середня довжина й час «вільного пробігу»

\begin{equation*}
	\bar{l} = \frac{1}{\sigma n/Z} = \frac{Z}{\sigma n}, \quad \bar{\tau} \approx \frac{\bar{l}}{v} = \frac{m^2\bar{v}^3}{\pi Z n e^4 L}.
\end{equation*}

Підставивши сюди $m\bar{v}^2 \approx 3kT$, отримаємо

\begin{equation}\label{eq:relaxation-time}
	\bar{\tau} \approx \frac{(3kT)^{3/2}\sqrt{m}}{\pi Z n e^4 L}.
\end{equation}

Звідси для питомої провідності плазми знаходимо

\begin{equation}\label{eq:plasma-conductivity}
	\lambda \approx \frac{ne^2\bar{\tau}}{m} \approx \frac{(3kT)^{3/2}}{\pi Z e^2 L\sqrt{m}}.
\end{equation}

Звичайно, наведений висновок треба розглядати як грубу оцінку, а не як доказ формули~\eqref{eq:plasma-conductivity}. Однак можна дати й більш обґрунтований висновок цієї формули. Підставляючи числові значення й виражаючи енергетичну температуру $\Theta = kT$ в електронвольтах, приходимо до оціночної формули для $\lambda$ (в обернених секундах):

\begin{equation*}
	\lambda \approx \frac{10^{13}}{Z}\Theta^{3/2},
\end{equation*}

причому ми прийняли $L \approx 15$. Провідність плазми росте пропорційно абсолютній температурі в степені $3/2$. У гарячій плазмі провідність стає дуже високою. Так, при $\Theta = 10$~кеВ для дейтерієвої плазми $\lambda \approx 10^{19}$~с$^{-1}$, тобто більше, ніж у міді ($5\cdot10^{17}$~с$^{-1}$). Ще швидше росте з температурою теплопровідність плазми, а саме пропорційно $\Theta^{5/2}$, так як для плазми, очевидно, повинен бути справедливий закон Відемана--Франца.

Велика різниця в масах електронів і іонів плазми робить можливим у плазмі існування таких квазірівноважних станів, які у відомому наближенні можуть бути охарактеризовані \emphx{двома температурами}. Дійсно, нехай початковий розподіл швидкостей електронів і іонів плазми ізотропний, але не максвеллівський. При зіткненні електрона з іншим електроном вони обмінюються енергією, величина якої порядка початкової енергії самих електронів. Тому час встановлення розподілу електронів за енергіями (тобто максвеллівського розподілу) із-за зіткнень між ними можна оцінити за формулою~\eqref{eq:relaxation-time}. (Це стане ясно, якщо в ній масу електрона $m$ замінити приведеною масою $m/2$.) Це час, який називається \emphx{електронним часом релаксації} $\bar{\tau}_e$, пропорційно квадратному кореню з маси електрона: $\bar{\tau}_e \sim \sqrt{m_e}$. Точно так же визначається \emphx{іонний час релаксації}, за яке встигає встановлюватися розподіл за енергіями між іонами із-за зіткнень між ними: $\bar{\tau}_i \sim \sqrt{m_i}$. Не то буде при взаємодії електронів з іонами. Тут за формулою~\eqref{eq:relaxation-time} получається час порядка $\sqrt{m_e}$. Однак при кожному зіткненні швидка частинка передає повільній лише незначну долю своєї енергії. В середньому ця доля порядка $m_e/m_i$ від першопочаткової енергії швидкої частинки (див.~т.~I, \S~28, задача~9). Для вирівнювання енергій потрібно релаксаційний час $\bar{\tau}_{ie}$, в $m_i/m_e$ раз більший, ніж $\bar{\tau}_e$. Таким чином,

\begin{equation}\label{eq:relaxation-times-ratio}
	\bar{\tau}_e : \bar{\tau}_i : \bar{\tau}_{ie} \sim 1 : \sqrt{\frac{m_i}{m_e}} : \frac{m_i}{m_e}.
\end{equation}

Звідси випливає, що $\bar{\tau}_e \ll \bar{\tau}_i \ll \bar{\tau}_{ie}$. Якщо плазму надати самій собі, то спочатку встановиться максвеллівський розподіл швидкостей електронів, потім --- іонів. Виникне квазірівноважний стан,
в якому електрони будуть мати температуру $T_e$, а іони --- температуру $T_i$. Взагалі кажучи, $T_e \neq T_i$. В цьому випадку плазму називають \emphx{неізотермічною} або \emphx{двотемпературною}. Потім у результаті обміну енергіями між електронами й іонами встановиться максвеллівський розподіл для всієї плазми, який характеризується загальною температурою електронів та іонів (\emphx{ізотермічна плазма}).

Коли плазма знаходиться в електричному полі, то в ній починає текти електричний ток і виділятися джоулеве тепло. При цьому енергію від поля отримують майже виключно електрони, як найбільш рухливі частинки. Іони нагріваються головним чином за рахунок енергії, яку вони отримують від «гарячих» електронів при кулонівських взаємодіях з ними. Так як останній процес відбувається порівняно повільно, то \emphx{температура електронів у плазмі виявляється вищою за температуру іонів}. Різниця між ними може бути вельми значною. Так, у позитивному стовпі тліючого розряду при тисках порядка 0,1 міліметрів ртутного стовпа температура електронів може досягати 50~000~°C і вище, тоді як температура іонів не перевищує декількох сотень градусів.

Основний практичний інтерес, який представляє фізика плазми, пов'язаний з рішенням проблеми \emphx{керованого термоядерного синтезу}. Для того щоб у речовині почалися достатньо інтенсивні термоядерні реакції, його необхідно нагріти до температури в декілька кілоелектронвольт або десятків кілоелектронвольт, а при таких температурах всяка речовина знаходиться в стані плазми. Найбільш перспективними «робочими речовинами» для термоядерного реактора являються \emphx{ізотопи водороду}: \emphx{дейтерій} (D) і \emphx{тритій} (T). Термоядерну реакцію синтезу легше отримати не в чистому дейтерії, а в його суміші з тритієм. Повне кількість дейтерію в океанах $4\cdot10^{13}$ тонн, що еквівалентно енергії $10^{20}$~кВт$\cdot$рік (повна споживана на всьому земному шарі потужність становить $10^{10}$~кВт). Тритій, як сильно радіоактивний елемент, в природних умовах не зустрічається, а получається штучно. В будущих термоядерних реакторах расход тритія должен с ізбитком пополняться воспроизводством (регенерацією) його в результаті опромінення $^6$Li нейтронами, получающимися в самих термоядерних реакторах.

Так як термоядерні реакції повинні відбуватися порівняно плавно й повільно, то виникає необхідність достатньо тривалого удержання гарячої плазми в обмеженому об'ємі робочої камери й ізоляції її від стінок цієї камери. Для цього пропонується використовувати \emphx{магнітну термоізоляцію}, тобто поміщати плазму в сильне магнітне поле, яке перешкоджає іонам і електронам переміщатися в поперечному напрямку й уходити на стінки камери.

Необхідне вимога, якому повинен задовольняти всякий термоядерний реактор, полягає в тому, щоб енергія, яка виділяється в ядерних реакціях, з ізбитком компенсувала затрати енергії від
зовнішніх джерел. Основними джерелами втрат енергії являються \emphx{гальмівне випромінювання} електронів при кулонівських зіткненнях останніх, а також \emphx{магнітогальмівне} (\emphx{циклотронне} або \emphx{бетатронне}) випромінювання, яке виникає внаслідок прискореного руху електронів у магнітному полі. Для самопідтримуючих термоядерних реакцій потрібно нагріти плазму до деякої «критичної» температури (50~кеВ для реактора на чистому дейтерії та 10~кеВ для реактора на рівнокомпонентній суміші дейтерію з тритієм). При цьому, як показує розрахунок, повинен виконуватися так званий критерій Лоусона: $n\tau > 10^{16}$~с/см$^3$ для реактора D--D і $n\tau > 10^{14}$~с/см$^3$ для реактора D--T. Тут $n$ --- концентрація іонів плазми (одного знака), а $\tau$ --- середній час удержання плазми.

Основна трудність, що стоїть на шляху створення керованого термоядерного синтезу, пов'язана з отриманням \emphx{спокійної}, або \emphx{стійкої}, \emphx{плазми}. Справа в тому, що із-за далекодіючого характеру кулонівських сил у плазмі відбуваються різні колективні процеси, наприклад мимоволі виникаючі \emphx{шуми й коливання}, які роблять плазму \emphx{нестійкою}. Основні зусилля при рішенні проблеми керованого термоядерного синтезу й направлені на подавлення цих нестійкостей.

\section*{Задачі}
\phantomsection
\addcontentsline{toc}{section}{Задачі}

\begin{problem}\label{prb:plasma-field}
	Напівпростір справа від нескінченної зарядженої площини $AB$ (рис.~\ref{pic:plasma-field}) заповнено повністю іонізованою плазмою, яка складається з однозарядних позитивних іонів і електронів. Знайти електричне поле й розподіл зарядів у плазмі на будь-яких відстанях від площини $AB$, якщо поверхнева густина зовнішніх зарядів на площині $AB$ рівна $\sigma$.
\end{problem}

\emphx{Розв'язання.} В плазмі поблизу площини $AB$ утворюється надлишок іонів, заряджених різнойменно з $\sigma$, і недостаток іонів, заряджених однойменно. В плазмі виникає електричне поле $\vect{E}$, перпендикулярне до площини $AB$. Направимо координатну ось $X$ усередину плазми перпендикулярно до цієї площини. Із міркувань симетрії ясно, що всі величини будуть функціями тільки координати $x$. Отже,

\begin{equation}\label{eq:poisson-1d}
	\div\vect{E} = \frac{dE}{dx} = 4\pi\rho.
\end{equation}

Густина вільних електричних зарядів $\rho$ може бути представлена у вигляді $\rho = (n^+ - n^-)e$, де $n^+$ --- концентрація позитивних, а $n^-$ --- негативних іонів. Вдалині від границі $AB$ (в нескінченності) плазма \emphx{електрично нейтральна}, там $n^+ = n^- \equiv n$. Поблизу границі концентрації $n^+$ і $n^-$ різні. Якщо $\varphi$ --- електричний потенціал поля $\vect{E}$, то згідно з формулою Больцмана

\begin{equation}\label{eq:boltzmann}
	n^+ = n\exp\left(-\frac{e\varphi}{kT}\right), \quad n^- = n\exp\left(+\frac{e\varphi}{kT}\right).
\end{equation}

При $x = +\infty$ потенціал $\varphi$ рівен нулю, а тому $n^+ = n^- = n$. Отже,

\begin{equation*}
	n^+ - n^- = n\left[\exp\left(-\frac{e\varphi}{kT}\right) - \exp\left(+\frac{e\varphi}{kT}\right)\right] = -2n\sh\frac{e\varphi}{kT}.
\end{equation*}

\begin{figure}[h!]\centering
	% \includegraphics[width=0.4\linewidth]{plasma-field}
	\caption{До задачі про поле в плазмі поблизу зарядженої площини}
	\label{pic:plasma-field}
\end{figure}

Отже,

\begin{equation*}
	\frac{dE}{dx} = -8\pi ne \sh\frac{e\varphi}{kT}.
\end{equation*}

Множачи обидві частини останнього рівняння на $2Edx = -2d\varphi$, отримаємо

\begin{equation*}
	dE^2 = 16\pi ne \sh\frac{e\varphi}{kT} d\varphi.
\end{equation*}

Просте інтегрування дає

\begin{equation*}
	E^2 = 16\pi nkT \ch\frac{e\varphi}{kT} + C.
\end{equation*}

Сталу $C$ треба визначити з умови, щоб у нескінченності (де $\varphi = 0$) поле $E$ оберталося в нуль. Це дає

\begin{equation*}
	E^2 = 16\pi nkT\left(\ch\frac{e\varphi}{kT} - 1\right) = 32\pi nkT \sh^2\frac{e\varphi}{kT}.
\end{equation*}

При добуванні квадратного кореня перед ним треба взяти знак плюс. Це випливає з того, що коли поле $E$ позитивно, тобто напрямлено вправо, то й потенціал $\varphi$ також позитивний. Якщо ж поле від'ємне, тобто напрямлено вліво, то й потенціал $\varphi$ від'ємний. Отже,

\begin{equation}\label{eq:plasma-field}
	E = -\frac{d\varphi}{dx} = 4\sqrt{2\pi nkT}\sh\frac{e\varphi}{kT}.
\end{equation}

Після інтегрування отримаємо

\begin{equation*}
	\th\frac{e\varphi}{kT} = C_1 e^{-x/D},
\end{equation*}

де $D$ --- стала:

\begin{equation}\label{eq:debye-radius-exact}
	D = \sqrt{\frac{kT}{8\pi ne^2}}.
\end{equation}

Ця формула й дає точний вираз дебаївського радіуса, введеного вище в тексті цього параграфа.

Сталу інтегрування $C_1$ можна виразити через потенціал $\varphi_0$ на границі плазми. Таким шляхом отримаємо

\begin{equation}\label{eq:potential-distribution}
	\th\frac{e\varphi}{kT} = \th\frac{e\varphi_0}{kT} e^{-x/D}.
\end{equation}

Самий потенціал $\varphi$ легко зв'язати з напруженістю зовнішнього електричного поля $E_0 = 2\pi\sigma$. При $x = 0$ повинно бути $E = E_0$, а тому

\begin{equation}\label{eq:field-at-boundary}
	E_0 = 4\sqrt{2\pi nkT}\sh\frac{e\varphi_0}{kT}.
\end{equation}

Таким чином, при віддаленні від границі плазми величина $\th(e\varphi/kT)$ експоненціально убуває. На протязі дебаївського радіуса вона убуває в $e$ раз. Цим визначається й закон убування поля $E$ з координатою $x$. Можна сказати, що за порядком величини дебаївський радіус $D$ визначає глибину, на яку електричне поле проникає в плазму.

\begin{problem}\label{prb:point-charge-plasma}
	Розв'язати попередню задачу для точкового заряду $q$, оточеного з усіх сторін плазмою. Обмежитися відстанями від заряду $r$, на яких виконується умова $|e\varphi| \ll kT$.
\end{problem}

\emphx{Розв'язання.} При виконанні умови $|e\varphi| \ll kT$ вирази~\eqref{eq:boltzmann} можна розкласти в ряди й обмежитися членами першої степені. Підставляючи отримані вирази в рівняння Пуассона $\Delta\varphi = -4\pi e(n^+ - n^-)$,
придамо йому вид
\begin{equation}\label{eq:screened-poisson}
	\Delta\varphi + \frac{\varphi}{D^2} = 0,
\end{equation}

де $D$ визначається попереднім виразом~\eqref{eq:debye-radius-exact}. Зважаючи на сферичну симетрію

\begin{equation*}
	\frac{1}{r^2}\frac{d}{dr}\left(r^2\frac{d\varphi}{dr}\right) + \frac{\varphi}{D^2} = 0.
\end{equation*}

Загальний розв'язок цього рівняння

\begin{equation}\label{eq:general-solution}
	\varphi = \frac{1}{r}\left(Ce^{-r/D} + C'e^{r/D}\right).
\end{equation}

Стала інтегрування $C'$ повинна обертатися в нуль, щоб задовольнити умову $\varphi = 0$ при $r = \infty$. Для визначення сталої $C$ допустимо, що умова $|e\varphi| \ll kT$ виконується вже при $r \ll D$. Тоді на таких відстанях ще применимо рішення~\eqref{eq:general-solution}. Но тоді воно переходить у кулонівський потенціал $\varphi = q/r$. Отже, повинно бути $C = q$, тобто

\begin{equation}\label{eq:debye-potential}
	\varphi = \frac{q}{r}e^{-r/D}.
\end{equation}

Цей потенціал називається \emphx{дебаївським потенціалом}. Формула~\eqref{eq:debye-potential} показує, що дебаївський радіус за порядком величини визначає відстань від заряду $q$, на якому кулонівське поле цього заряду екранується протилежно зарядженими іонами плазми. Можна також сказати, що дія кулонівського поля заряду $q$ простирається на відстань порядка дебаївського радіуса $D$, а на більших відстанях практично не має місця.

\begin{problem}\label{prb:surface-layer-thickness}
	В електростатиці вважається, що в стані рівноваги електрику розподіляється по поверхні провідника. Фактично воно розподіляється не по математичній поверхні, а всередині поверхневого шару кінцевої товщини $l$. Оцінити цю товщину.
\end{problem}

\emphx{Відповідь.} Величина $l$ порядка дебаївського радіуса $D$, якщо вираз для $D$ змінити з урахуванням розподілу Фермі для електронів металу. Для цього величину $kT$ треба замінити енергією Фермі~\eqref{eq:fermi-energy}. Зробивши це й опускаючи числові множники порядка одиниці, отримаємо

\begin{equation*}
	l = \sqrt{a_1 n^{-1/3}},
\end{equation*}

де $a_1 = h^2/(4\pi^2 me) = 0{,}53\cdot10^{-8}$~см (радіус першої боровської орбіти в атомі водороду), а $n^{-1/3}$ --- середня відстань між вільними електронами металу.

